\documentclass[11pt]{article}

\usepackage[margin=1in]{geometry}
\usepackage{amsmath,amssymb,amsfonts}
\usepackage{booktabs}
\usepackage{graphicx}
\usepackage{xcolor}
\usepackage{algorithm}
\usepackage{algpseudocode}
\usepackage{hyperref}

\title{Operator Identification for Nonlinear PDEs: \\
       MILP and CCP Algorithms}
\author{}

\begin{document}
\maketitle

\section{Problem Statement}
Given a PDE library $\Theta \in \mathbb{R}^{n \times P}$ (columns are candidate
differential operators evaluated on $n$ space-time grid points) and a
time-derivative vector $y \in \mathbb{R}^n$, find a sparse coefficient vector
$\xi \in \mathbb{R}^P$ such that $y \approx \Theta \xi$.
The true PDE has only $k \ll P$ active terms.

\subsection{Notation}
\begin{itemize}
  \item $n$ -- number of space-time snapshots (e.g.\ $128 \times 100 = 12800$)
  \item $P$ -- library size (S=12, M=23, L=40, XL=59, $\ldots$, P196=196)
  \item $\xi$ -- coefficient vector to recover
  \item $\|\cdot\|_p$ -- standard $\ell_p$ norm
\end{itemize}

% ===================================================================
\section{MILP: Mixed-Integer Linear Programming (L0 Baseline)}
% ===================================================================

The MILP approach solves the $\ell_0$-regularized problem via
a Big-M formulation and a bisection-based Pareto sweep.

\subsection{Formulation}

For a fixed tolerance $\varepsilon > 0$, the cardinality minimization is:
\begin{equation}
\min_{\xi, z} \; \sum_{i=1}^P z_i
\quad \text{s.t.} \quad
\|y - \Theta \xi\|_1 \le \varepsilon, \qquad
| \xi_i | \le M \, z_i, \quad z_i \in \{0,1\}
\label{eq:milp}
\end{equation}
where $M = \max(2\,\|\xi_{\text{OLS}}\|_\infty, 1)$ is computed from
the column-normalized OLS solution.

The $\ell_1$ constraint is linearized by introducing slack variables
$s \in \mathbb{R}^n$:
\begin{align}
y - \Theta\xi &\le s, \qquad -(y - \Theta\xi) \le s, \qquad \sum_{j=1}^n s_j \le \varepsilon
\end{align}
The resulting MILP has $2P + n$ variables ($P$ continuous, $P$ binary, $n$ continuous slack).

\subsection{Bisection-Based Pareto Sweep}

Algorithm~\ref{alg:pareto} performs a geometric bisection on $\varepsilon$
to trace the Pareto frontier $k(\varepsilon)$ (sparsity vs tolerance).

\begin{algorithm}[ht]
\caption{Bisection-based Pareto sweep}
\label{alg:pareto}
\begin{algorithmic}[1]
\State \textbf{Input:} $\Theta, y$, noise floor $\eta$, factors $f_{\text{lo}}, f_{\text{hi}}$
\State $\varepsilon_{\text{lo}} \gets \eta \cdot f_{\text{lo}}$
\State $\varepsilon_{\text{hi}} \gets \eta \cdot f_{\text{hi}}$
\State $k_{\text{lo}}, \xi_{\text{lo}} \gets \text{Solve}(\varepsilon_{\text{lo}})$
\State $k_{\text{hi}}, \xi_{\text{hi}} \gets \text{Solve}(\varepsilon_{\text{hi}})$
\State $K \gets \{k_{\text{lo}}: \varepsilon_{\text{lo}}, \; k_{\text{hi}}: \varepsilon_{\text{hi}}\}$
\For{$k \gets k_{\text{hi}}+1$ \textbf{to} $k_{\text{lo}}-1$}
  \State $\varepsilon_L \gets K(k+1)$ \Comment{tight -- not yet feasible at $k$}
  \State $\varepsilon_R \gets K(k-1)$ \Comment{loose -- feasible at $k$}
  \For{$t \gets 1$ \textbf{to} $N_{\varepsilon}$}
    \State $\varepsilon_m \gets \sqrt{\varepsilon_L \cdot \varepsilon_R}$
    \State $k_m, \_ \gets \text{Solve}(\varepsilon_m)$
    \If{$k_m \le k$} $\varepsilon_R \gets \varepsilon_m$ \Else $\varepsilon_L \gets \varepsilon_m$ \EndIf
  \EndFor
  \State $K(k) \gets \varepsilon_R$
\EndFor
\State Select plateau with smallest $\varepsilon$ lower bound (tightest tolerance)
\State \Return $\xi_{\text{selected}}$
\end{algorithmic}
\end{algorithm}

\subsection{Plateau Selection}

After the sweep, supports are grouped into runs of consecutive $\varepsilon$ values
with identical support. The \emph{non-trivial} plateau with the smallest
$\varepsilon$ lower bound is selected:
\begin{equation}
k^* = \arg\min_{\text{runs}} \; \varepsilon_{\text{lower}}(r), \qquad r \notin \{0\text{-support}\}
\end{equation}

\subsection{Adaptive Tightening}

After selecting a support, $\varepsilon_{\text{lo}}$ is repeatedly halved.
If the support remains unchanged for 3 consecutive halvings, the search stops.
If the support \emph{changes}, the new support replaces the old one and the
counter resets. This catches cases where the true support requires a tighter
tolerance than the initial sweep (e.g.\ Allen--Cahn with $\nu = 0.01$).

% ===================================================================
\section{CCP: Correlation-Cut Pursuit}
% ===================================================================

CCP takes a fundamentally different approach: instead of solving a global
MILP, it recursively partitions the feature set by spectral clustering and
performs local OLS voting between cluster pairs.

\subsection{Algorithm Overview}

\begin{enumerate}
  \item \textbf{Correlation graph:} $C_{ij} = |\Theta_i^\top \Theta_j| / n$,
    diagonal set to zero.
  \item \textbf{Spectral clustering:} Compute the graph Laplacian
    $L = D - C$, then its Fiedler vector (second eigenvector of $L$).
    Split along the Fiedler ordering recursively until each cluster has
    $\le$ \texttt{cluster\_size} features. Stop early if
    $\lambda_2 > \lambda_{\text{stop}}$ (the subgraph is well-connected).
  \item \textbf{Cross-group OLS voting:}\\
    For each pair of groups $(g_a, g_b)$:
    \begin{itemize}
      \item Concatenate columns into $T_{ab}$
      \item Column-normalize: $T_{ab} / \|T_{ab}\|_2$
      \item Solve $\min_{\xi} \|y - T_{ab}\xi\|_2^2$ (OLS)
      \item Vote for column $j$ if $|\xi_j| > \tau \cdot \max|\xi|$
        ($\tau = 10^{-3}$)
    \end{itemize}
  \item \textbf{Survivor selection:} Terms with votes $\ge \lceil n_g/2 \rceil$
    survive. If nothing pruned, tighten to $n_g - 1$ votes.
  \item \textbf{Iterate:} Re-cluster survivors; repeat until support
    $\le$ \texttt{cluster\_size} or no change.
  \item \textbf{Final OMP + OLS:} If $>6$ survivors, run OMP to pick 6;
    debias with OLS; prune coefficients below $10^{-4} \cdot \max|coef|$.
\end{enumerate}

\subsection{Key Design Decisions}

\textbf{Column normalization.} Raw OLS coefficients are dominated by column
magnitude (e.g.\ $u^3$ has much larger entries than $u$). Normalizing
to unit $\ell_2$ norm makes coefficients comparable. This was the key
fix that improved CCP from 10\% to 95\% on Burgers at $P=40$.

\textbf{Relative threshold.} The threshold $\tau = 10^{-3}$ matches the
Lasso pruning heuristic. It filters out terms whose normalized OLS coefficient
is below 0.1\% of the dominant term. The original threshold $10^{-14}$
let every non-zero term through, causing false positives at large $P$.

\textbf{Fiedler stop.} The adaptive stopping criterion
$\lambda_2 > \lambda_{\text{stop}}$ prevents over-cutting well-connected
subgraphs, avoiding degenerate small clusters that amplify noise.

% ===================================================================
\section{Computational Comparison}
% ===================================================================

\textbf{MILP (L0\_SCIP/L0\_CBC):}
\begin{itemize}
  \item $n_{\varepsilon} = 30$ bisection iterations per seed
  \item $\approx 30$--40 MILP solves per seed
  \item Each MILP: $2P + m$ variables ($P$ binary), $m = 500$ subsampled rows
  \item SCIP: 12K LP iterations per solve, 137 cutting planes at root node,
    $\sim$46s/solve on $P=40$ (Allen--Cahn)
  \item Total: 15--30 minutes per seed at $P=40$
\end{itemize}

\textbf{CCP:}
\begin{itemize}
  \item Spectral clustering via LAPACK \texttt{dsyevx} with
    \texttt{subset\_by\_index=[1,1]} (Fiedler vector only)
  \item $\binom{n_g}{2}$ OLS solves per round ($n_g = P / \text{cluster\_size}$)
  \item Each OLS: $\le 2\cdot\text{cluster\_size}$ columns $\times$ 12800 rows
  \item Total: 0.5--10s per seed (scaling roughly $O(P^2)$ from the Fiedler
    eigenvalue)
\end{itemize}

% ===================================================================
\section{Benchmark Results}
% ===================================================================

\subsection{CCP Across Library Sizes ($P=12$ to $P=196$)}
\label{sec:ccp_scaling}

Table~\ref{tab:ccp_scaling} shows CCP (default \texttt{cluster\_size=8}) across
16 library sizes, 7 PDEs, 20 seeds each. Values are Jaccard $J=1$ counts.

\begin{table}[ht]
\centering
\caption{CCP Jaccard $J=1$ across library sizes ($n=20$ seeds per cell)}
\label{tab:ccp_scaling}
\footnotesize
\begin{tabular}{lcccccccccccccccc}
\toprule
PDE   & 12 & 23 & 40 & 59 & 45 & 55 & 75 & 81 & 89 & 96 & 126 & 133 & 146 & 154 & 174 & 196 \\
\midrule
KdV     & 20 & 19 & 19 & 20 & 20 & 19 & 19 & 19 & 20 & 20 & 18 & 16 & 19 & 18 & 14 & 16 \\
Burgers & 20 & 20 & 19 & 19 & 20 & 19 & 17 & 16 & 15 & 16 & 13 & 9  & 11 & 15 & 9  & 6  \\
AC      & -- & -- & 20 & 20 & 20 & 20 & 19 & 18 & 19 & 18 & 17 & 19 & 18 & 19 & 17 & 17 \\
KS      & -- & 18 & 16 & 8  & 9  & 8  & 8  & 8  & 5  & 9  & 4  & 4  & 2  & 2  & 3  & 1  \\
FKPP    & 20 & 20 & 14 & 9  & 11 & 15 & 12 & 17 & 10 & 10 & 7  & 7  & 3  & 5  & 4  & 3  \\
KdVB    & 20 & 20 & 19 & 19 & 19 & 17 & 11 & 10 & 8  & 11 & 7  & 5  & 5  & 4  & 4  & 1  \\
FHN     & -- & -- & 18 & 17 & 19 & 18 & 13 & 9  & 14 & 9  & 10 & 7  & 9  & 4  & 6  & 6  \\
\bottomrule
\end{tabular}
\end{table}

KdV and Burgers maintain $J=1$ up to $P\sim150$. AC and FHN are near-perfect
through $P\sim170$. KS and FKPP degrade fastest due to high feature
correlation. KdVB is fragile above $P\sim80$ (3 true terms across clusters).

\subsection{MILP vs CCP at $P=40$ (L Library)}

Table~\ref{tab:milp_vs_ccp} compares L0\_SCIP, L0\_CBC, and CCP on the
standard L library ($P=40$), 7 PDEs, 20 seeds.

\begin{table}[ht]
\centering
\caption{MILP vs CCP at $P=40$ (Jaccard $J=1$ / 20 seeds)}
\label{tab:milp_vs_ccp}
\begin{tabular}{lcccc}
\toprule
PDE   & L0\_SCIP & L0\_CBC & CCP \\
\midrule
KdV     & 20/20 & 16/20 & 19/20 \\
Burgers & 20/20 & 20/20 & 19/20 \\
AC      & --    & --    & 20/20 \\
KS      & --    & --    & 16/20 \\
FKPP    & 20/20 & 19/20 & 14/20 \\
KdVB    & 20/20 & 9/20  & 19/20 \\
FHN     & 20/20 & --    & 18/20 \\
\bottomrule
\end{tabular}
\end{table}

\textbf{Notes:}
\begin{itemize}
  \item SCIP and CBC results from $n_{\varepsilon}=30$ bisection runs.
  \item Dashes ($--$) indicate CoinError solver crash (CBC on AC, KS, FHN)
    or solver not installed (SCIP on AC, KS).
  \item KdVB L CBC: 11/20 CoinError failures (pre-marked); SCIP provides
    $J=1$ for all seeds.
  \item SCIP solves take 15--30 min/seed at $P=40$; CCP takes 0.5--2s/seed.
  \item CCP outperforms MILP on KS, AC; MILP edges on FKPP, KdVB.
\end{itemize}

\subsection{Solver Comparison: SCIP vs CBC}

Both solve the identical Big-M MILP formulation (Eq.~\ref{eq:milp}) with
the same subsample ($m=500$, identical random seed). Differences arise from
solver internals:

\begin{itemize}
  \item \textbf{Presolve:} SCIP tightens 1602 variable bounds in 5 rounds;
    CBC's presolve is less aggressive.
  \item \textbf{Cutting planes:} SCIP generates 137 cuts in the root node,
    closing the gap from 621\% to 114\% before branching. CBC generates fewer,
    weaker cuts.
  \item \textbf{Heuristics:} SCIP's ``integrality shift'' finds feasible
    solutions faster (primal bound 2.0 found at root). CBC relies more
    on branching.
  \item \textbf{Branching:} SCIP explores 1--2 nodes (total 2 runs after
    restart with 28 global fixings). CBC may explore deeper trees.
  \item \textbf{Robustness:} CBC crashes with CoinError on ill-conditioned
    instances (AC, KS, FHN at $P=40$). SCIP handles all 7 PDEs without
    crashing.
\end{itemize}

\subsection{Solver Runtime Breakdown}

For a single Allen--Cahn seed ($P=40$, $n=12800$):

\begin{center}
\begin{tabular}{lccc}
\toprule
Phase & SCIP (s) & CBC (s) & CCP (ms) \\
\midrule
One MILP solve ($\varepsilon=100$) & 46.3 & 2--25 & -- \\
$\times$ 30 bisection steps & 1389 & 60--750 & -- \\
Total (one seed) & $\sim$25 min & $\sim$2--12 min & $\sim$500 ms \\
\bottomrule
\end{tabular}
\end{center}

CCP's cross-group OLS voting ($\sim$28 pairs $\times$ OLS on $\le$16 columns)
replaces the global MILP with $\sim$28 cheap LS solves, achieving a
$1000\times$ speedup.

% ===================================================================
\section{Discussion}
% ===================================================================

\subsection{Why CCP Outperforms MILP on KS/AC}

\textbf{Allen--Cahn ($\nu=0.01$):} The true support has coefficients
$\{u^3: 382, u: 136, u_{xx}: 21\}$. The Big-M $M \approx 764$ treats all
three identically. The MILP can drop $u_{xx}$ (coefficient 21) because
the $\varepsilon$-tolerance absorbs its residual. More bisection iterations
($n_{\varepsilon}=30$) help but don't fully solve this -- the adaptive
tightening added in this work addresses it by halving $\varepsilon_{\text{lo}}$
until the support stabilizes.

CCP succeeds because $u_{xx}$ is clustered into a small group (it has low
correlation with $u$ and $u^3$) where OLS easily identifies it.

\textbf{Kuramoto--Sivashinsky:} The library contains many terms correlated
with $u_{xxxx}$ ($u^2 u_{xxxx}, u^3u_{xxxx}, \ldots$). The OLS coefficients
show no clear separation between true terms ($u_{xxxx}: 3110$) and
spurious ($u^2 u_{xxxx}: 479$). The MILP's Pareto front is flat (many
sparsity levels achieve similar $\varepsilon$), making plateau selection
unreliable.

CCP's clustering groups correlated terms together; the OLS voting then
selects the correct representative from each group.

\subsection{When to Use MILP}

MILP (SCIP) is useful when:
\begin{itemize}
  \item The library is small ($P \le 23$) -- solves are fast (2--5s)
  \item The coefficient magnitudes are well-separated (KdV, Burgers,
    KdVB)
  \item A provable optimality certificate is needed
  \item The CCP voting fails due to wide term spread across clusters
\end{itemize}

For large libraries ($P \ge 40$), CCP is the preferred method: faster,
more robust, and equally or more accurate on most PDEs.

\end{document}
